\section{Experimental Setup}
\label{sec:setup}

\subsection{Arms}
\label{sec:arms}

All six arms share one implementation, one schedule, one budget and one seed
list; they differ only in how the mixing operator of \eqref{eq:master} is
constructed. Table~\ref{tab:arms} states the differences exactly.

\begin{table}[t]
\caption{The six arms. All other components---schedule, elite memory,
mutation, boundary handling, budget, seeds---are identical. Degree is fixed
except in the last two arms, which redraw it every $12$ iterations.}
\label{tab:arms}
\centering
\small
\setlength{\tabcolsep}{4pt}
\begin{tabular}{@{}llll@{}}
\toprule
Arm & Degree & Weights & Role \\
\midrule
\textsc{Social}     & $5$      & betweenness & published \\
\permcontrol{}      & $5$      & permuted    & control 1 \\
\textsc{Unif}$_4$   & $5$      & uniform     & ablation \\
\textsc{Fix}$_{16}$ & $16$     & uniform     & fixed density \\
\polcontrol{}       & random   & uniform     & control 2 \\
\textsc{Ads}        & UCB      & uniform     & adaptive \\
\bottomrule
\end{tabular}
\end{table}

\permcontrol{} draws a fresh permutation whenever the graph changes.
\polcontrol{} and \textsc{Ads} both re-draw the neighbourhood degree from
$\{4,8,16,32,N{-}1\}$ every $12$ iterations and rebuild the graph at that
degree, carrying positions and fitness over; they differ only in whether the
choice uses the observed improvement (\textsc{Ads}, sliding-window UCB with
window $40$ and $c=0.35$) or ignores it (\polcontrol{}, uniform draw). Both
use uniform mixing, since Section~\ref{sec:results} shows the centrality
weighting to be uninformative; dropping it also removes the betweenness
computation, about $30\%$ of per-iteration runtime.

\subsection{Benchmarks and protocol}
\label{sec:benchmarks}

\emph{Classic suite.} The suite of Yao \emph{et
al.}~\cite{yao1999evolutionary}, used because it is the one on which the method
under test was originally reported. The control comparison uses ten of its $23$
functions---\textsc{Sphere}, \textsc{Schwefel}~1.2, \textsc{Rosenbrock},
\textsc{Step}, \textsc{Quartic}, \textsc{Schwefel}~2.26, \textsc{Rastrigin},
\textsc{Ackley}, \textsc{Griewank} and \textsc{Penalized}---one or two from
each of its unimodal, discontinuous, multimodal and penalised groups, chosen
before any result was seen. The remaining thirteen are fixed-dimension
problems, which the density sweep of Section~\ref{sec:res:density} runs at
their native dimensions but which are ill-suited to a paired contrast at
$D=30$ (Section~\ref{sec:corrections}). We run the ten at two budgets,
$2\times10^{4}$ and $1.05\times10^{5}$ evaluations, the latter matching the
original report.

\emph{CEC2017.} The competition suite~\cite{awad2016cec2017} with the
prescribed protocol: search domain $[-100,100]^{D}$, budget $10^{4}D$
evaluations, and error $f(\mathbf{x})-f^{*}$ as the reported outcome. We use
$D\in\{10,30,50\}$. Two functions are absent: $F_{2}$, which the organisers
advise against for unstable behaviour in higher dimensions, and $F_{30}$,
which the reference implementation we use does not provide---leaving $28$
functions.

\emph{Repetitions.} The primary contrast---the published method against its
permutation control---is run at $D=30$ for the $51$ repetitions the CEC2017
protocol prescribes, and every number we report for that contrast at $D=30$
uses all $51$. The other four arms, and all arms at $D\in\{10,50\}$, are run
for $30$. Comparisons involving those arms, including the six-arm table and
rank ordering in the supplementary material, therefore use the common $30$-run
subset of all six; mixing run counts within one comparison would break the
pairing the statistics rely on.

Population size is $N=60$ throughout. Seeds are $0,1,\dots$ shared across
arms, so every comparison is paired: an arm and its control start from the
same population and consume the same random stream.

\subsection{Corrections, implementation and cost}
\label{sec:corrections}

Three defects in the published implementation affect measurement and are
corrected here: the fixed-dimension functions $F_{14}$--$F_{23}$ are
evaluated at $D=30$ although their implementations use only the first two to
six coordinates, leaving up to $28$ dimensions inert; neighbourhood size is
set per function category, which is per-instance tuning; and the rewiring
routine catches an exception class that does not cover the one
\texttt{double\_edge\_swap} raises, so dense graphs crash the optimizer.

The supplementary material states each correction, reports the equivalence
check of our vectorized update against the reference path
($8.9\times10^{-16}$ per diffusion step, identical run distributions), and
gives the hardware, the measured parallel throughput and the twentyfold
spread of per-run cost across the suite.
